\documentclass[12pt,a4paper]{ctexart}
\usepackage[a4paper,margin=2.4cm]{geometry}
\usepackage{fontspec}
\usepackage{enumitem}
\usepackage[hidelinks]{hyperref}
\usepackage{xcolor}
\usepackage{graphicx}

\setmainfont{Times New Roman}
\setCJKmainfont{SimSun}[BoldFont=SimHei]
\setCJKsansfont{SimHei}
\linespread{1.18}
\emergencystretch=3em
\sloppy
\setlength{\parindent}{2em}
\setlength{\parskip}{0.35em}
\setlist{leftmargin=2.2em,itemsep=0.35em,topsep=0.3em}
\pagestyle{plain}

\ctexset{
  section = {
    format = \large\bfseries\sffamily,
    beforeskip = 1.1em,
    afterskip = 0.7em
  }
}

\begin{document}

\begin{center}
\includegraphics[width=2.8cm,height=2.8cm,keepaspectratio]{assets/profile-avatar-clean.png}\\[0.5em]
{\LARGE\bfseries 程振}
\end{center}

\section*{一、个人基本情况}

个人主页：\href{https://upc20zcheng.github.io/}{https://upc20zcheng.github.io/}

基本信息：程振，男，汉族，安徽利辛人，讲师，硕士生导师，毕业于中国石油大学（华东）管理科学与工程专业并获得管理学博士学位。

研究方向：绿色智能财会、会计数据标准增信、算电协同与智能经济新形态

项目主持及参与：主持安徽省社科基金青年项目1项、厅市级项目2项、阜阳师范大学青年人才基金重点项目1项。参与四项国家、省级与教育部课题，以及一项横向课题。

论文发表：近五年在以第一作者/通讯作者在\textit{Energy}、\textit{Journal of Cleaner Production}、中国人口·资源与环境、产业经济研究、商业研究、经济问题探索、江淮时报、技术经济与管理研究等国内外学术期刊（中科院1区/国家自然科学基金委认定的B类/FMS高质量中文期刊/CSSCI）上发表学术论文10余篇。

学术兼职：\textit{International Review of Economics \& Finance}、\textit{International Review of Financial Analysis}、\textit{Environmental Engineering and Management Journal}、\textit{Frontiers in Public Health}、\textit{PLOS ONE}等多个期刊审稿人。

在授课程：《管理定量分析：方法与技术》《管理会计学》（本科生）、《统计分析方法》《财务管理理论与实务》（研究生）

\section*{二、近五年代表性论文}

\begin{enumerate}[label={[\arabic*]}]
\item 第一作者. Energy use rights trading and carbon emissions[J]. \textit{Energy}, 2025, 315: 134360. （中科院SCI 一区）
\item 第一作者. Can environmental information disclosure promote urban carbon emission reduction? Quasi-experimental evidence from China[J]. \textit{Journal of Cleaner Production}, 2025, 489: 144698.（中科院SCI 一区）
\item 第一作者.新质生产力视域下安徽构建“政产学研金服用”七位一体创新生态实现非对称赶超研究[N].江淮时报,2025-08-29(006).（省级报刊·理论版）
\item 通讯作者.中国环境保护税法对企业劳动雇佣的影响[J].中国人口·资源与环境,2023,33(01):61-73.（CSSCI、国家自然科学基金委认定的B类期刊、FMS T2）
\item 通讯作者.国家低碳战略提高了企业全要素生产率吗?——基于低碳城市试点的准自然实验[J].产业经济研究,2021,(06):101-115.（CSSCI、FMS高质量中文期刊 T2）
\item 通讯作者.黄河流域数字经济、生态保护与高质量发展时空耦合及其驱动因素[J].经济问题探索,2022,(08):135-148.（CSSCI）
\item 通讯作者.低碳城市建设能成为企业数字化转型的新动能吗?——来自低碳城市试点的经验证据[J].商业研究,2024,(03):80-91.（CSSCI）
\item 通讯作者.环境规制对我国城市碳排放的影响——门槛特征与空间效应[J].技术经济与管理研究,2024,(11):67-72.（北大中文核心）
\item 第一作者.绿色信贷、企业风险承担与企业绿色创新——“波特效应”的中国证据[J].技术经济与管理研究,2022,(10):68-74.（北大中文核心）
\end{enumerate}

\section*{三、近五年主持/参与的主要教科研项目}

\begin{enumerate}[label=\arabic*.]
\item 安徽省社科规划青年项目（AHSKYQ2024D022，主持）。
\item 安徽省财政科研课题（AHCZ2025LX12，主持）。
\item 阜阳市社科规划课题（FSK2025059，主持）。
\item 阜阳师范大学青年人才基金重点项目（，主持）。
\item 国家社科基金青年项目（21CJY035，参与）。
\item 国家社科基金一般项目（20BJY090，参与）。
\item 中国石油天然气股份有限公司长庆油田分公司项目（21YL2-FW-011，参与）。
\item 山东省社会科学规划金融研究专项（19CJRJ02，参与）。
\item 中央高校基本科研业务费专项（19CX04034B，参与）。
\end{enumerate}

\section*{四、获奖情况}

\begin{enumerate}[label=\arabic*.]
\item “第七届营销科学与创新国际高峰论坛（MSI 2023）暨《\textit{Internet Research}》专辑国际学术会议”，“优秀论文二等奖”。
\item The Competition on Energy and Climate Change (CECC), a special event of BRICS Youth Energy Summit \& International Energy Youth Conference，Worldwide Energy University Network, Excellence Award。
\end{enumerate}

\section*{五、联系方式}

电子邮箱：\href{mailto:202406038@fynu.edu.cn}{202406038@fynu.edu.cn}、\href{mailto:13237494718@163.com}{13237494718@163.com}

联系地址：安徽省阜阳市清河东路741号；邮政编码236041

\end{document}
